\subsection{Data Model}
\label{sec:data-model}

\begin{longtable}{|p{3cm}|p{0.68\textwidth}|}
\caption{Initial Breakly data model}\label{tab:data-model}\\
\hline
\textbf{Entity} & \textbf{Main attributes} \\
\hline
\endfirsthead
\multicolumn{2}{r}{\textit{Continued from previous page}}\\[-2pt]\hline
\textbf{Entity} & \textbf{Main attributes} \\
\hline
\endhead
\multicolumn{2}{r}{\textit{Continue on next page}}\\
\endfoot
\endlastfoot
UserPreferences & userId, defaultDuration, reminderEnabled, reminderSchedule,
accessibilityPreferences \\
\hline
Activity & id, title, description, estimatedDuration, category, instructions \\
\hline
BreakSession & id, activityId, startedAt, endedAt, plannedDuration, status \\
\hline
Reminder & id, schedule, enabled, notificationPermissionState \\
\hline
\end{longtable}

The ERD will define that a user has one preference set and many sessions;
each session references one activity, while reminders belong to the user's
notification configuration. Timestamps are stored in a consistent format,
and deleted or reset data follows an explicit retention policy.

\subsection{Local Storage Strategy}
\label{sec:local-storage}

Preferences and session history are persisted locally so the main experience
works offline. The timer derives remaining time from a stored start timestamp
and planned duration rather than relying only on elapsed foreground ticks.
Repositories hide the storage implementation from the domain layer. The
report will document schema migrations, reset behaviour, and corruption
recovery.

\subsection{API Contract}
\label{sec:api-contract}

Optional synchronisation endpoints will follow a versioned JSON contract:
\begin{itemize}
\item \texttt{GET /v1/activities} returns available activities and metadata.
\item \texttt{GET /v1/preferences} and \texttt{PUT /v1/preferences} read and
  update user preferences.
\item \texttt{POST /v1/sessions} records a completed session.
\item \texttt{GET /v1/sessions} returns paginated session history.
\end{itemize}

Successful responses will define explicit schemas, while errors use a
consistent object containing \texttt{code}, \texttt{message}, and optional
\texttt{details}. The contract must document 400 validation errors, 401
authentication errors, 404 missing resources, 409 conflicts, and 429/5xx
temporary failures, together with retry and offline fallback behaviour.
